\section{Related Works}

\paragraph{Fine-grained learning in low-data regimes.}
Fine-grained learning under limited supervision has been studied mainly for classification. The survey in \cite{tang2024awesome_fgfs} shows that prior work is largely restricted to image-level recognition, with little attention to dense prediction. Fine-grained semantic segmentation is more demanding, as it requires both recognition and precise localization. Only a small number of recent works, such as \cite{part_matching}, address related settings, mainly through part-level correspondence rather than category-level segmentation.

\paragraph{FungiTastic benchmark and prior solutions.}
FungiTastic \cite{Picek_2025_CVPR} has so far been explored mostly for classification. Recent FungiCLEF submissions, such as \cite{prototypical_networks, contrastive}, focus on long-tailed recognition through feature engineering, ensembling, and lightweight finetuning. These methods target classification and do not address fine-grained segmentation.

\paragraph{Prototype-based methods.}
Prototype-based learning is a standard approach for few-shot recognition \cite{proto1, proto2, proto3}, where classes are represented in an embedding space. Recent works also use multimodal prototypes to improve generalization under class imbalance \cite{multimodal}. Our method follows this idea, but combines prototype-based classification with class-agnostic segmentation in a single pipeline.

\paragraph{PCA whitening and feature preprocessing.}
PCA whitening is a common preprocessing step for distance-based inference, as it decorrelates features and normalizes variance across directions \cite{kessy2018optimal}. It is particularly useful when pretrained embeddings are dominated by nuisance variation rather than task-relevant structure \cite{anthony2026invisible, weng2022investigation}. This makes it a natural choice for prototype matching in our setting.